\subsection{0. Executive decision and theory-to-implementation bridge audit}\label{executive-decision-and-theory-to-implementation-bridge-audit}

\subsubsection{0.1 Central question}\label{central-question}

How can an agent preserve an authorized goal as a durable commitment across many execution cycles while still reconsidering whether the current method remains appropriate, and who has authority to continue, delegate, suspend, abandon, or change that method?

The chapter studies B4: \textbf{persistent intention, explicit reconsideration, and live method-switch authority}. It is not a generic chapter about planning, reflection, memory, goal setting, or prompt injection. Those topics appear only where they illuminate the B4 control relation.

\subsubsection{0.2 The B4 claim in Fan \& Lan's terms}\label{the-b4-claim-in-fan-lans-terms}

Fan \& Lan propose that a persistent intention record should retain the committed goal, rationale, plan, continuation condition, reconsideration triggers, authority, and termination status. A governor should decide whether to continue, delegate, suspend, abandon, or change method. The closest reviewed baselines include hybrid and self-aware BDI-LLM systems, CCA, Devil's Advocate, ADAS, and AFlow. B4 is weakened or rejected if commitments become stale, reconsideration oscillates, the agent stops prematurely, or method changes have no diagnostic value.

The chapter must distinguish what Fan \& Lan state from what the book designs. The paper proposes a residual control bundle and identifies evidence and failure conditions; it does not prescribe a complete production schema, transition table, authorization protocol, concurrency model, or universal reconsideration score.\FanLanCite{\S VI-C; Table VI; Figure 8; Appendix D, Table X; Appendix F, R5a--R5c and R6}

\subsubsection{0.3 Theory-to-practice assessment}\label{theory-to-practice-assessment}

\textbf{Assessment: a substantial implementation bridge is necessary.} The historical theories identify distinct state and control relationships, but a reader still needs to know how a modern runtime represents an intention, preserves it between model calls, detects a justified reconsideration trigger, authorizes a method change, and prevents that change from silently rewriting the goal.

The outline therefore requires the following implementation material:

\begin{enumerate}
\def\labelenumi{\arabic{enumi}.}
\tightlist
\item
  A reference architecture with a principal/task interface, immutable goal specification, persistent intention registry, belief/evidence store, execution monitor, reconsideration detector, diagnosis component, method registry, method-authority governor, authorization gate, executor, delegation boundary, and append-only decision ledger.
\item
  Typed contracts for \texttt{GoalSpec}, \texttt{IntentionRecord}, \texttt{MethodDescriptor}, \texttt{MethodAssignment}, \texttt{ReconsiderationEvent}, \texttt{GovernorDecision}, \texttt{DelegationContract}, and \texttt{AuthorityPolicy}.
\item
  A lifecycle state machine covering proposal, authorization, active execution, reconsideration, continuation, method switch, delegation, suspension, satisfaction, impossibility, and abandonment.
\item
  Explicit authority separation: the principal controls the goal; the governor may choose only among permitted methods; a model may propose but cannot commit; the monitor diagnoses; safety and authorization controls may veto; external side-effect permission remains a separate boundary.
\item
  Persistence and consistency rules: versioned intention records, append-only events, optimistic concurrency, deduplication, restart recovery, method-switch transactions, stale-evidence checks, and audit reconstruction.
\item
  Reconsideration controls: evidence-gated triggers, hysteresis, cooldowns, switch budgets, escalation, and a distinction between a new plan under the same method and an explicit method switch.
\item
  Delegation and suspension contracts, including scope, expiry, return conditions, and the requirement that delegated work does not acquire authority to amend the parent goal.
\item
  A worked Product X payment-status example, a vertical-slice implementation sequence, Java-like contracts/pseudocode, and executable tests.
\item
  An evaluation design faithful to Fan \& Lan Appendix D: required state trace, primary and safety/efficiency outcomes, baselines, three specified ablations, and falsification criteria.
\end{enumerate}

These are \textbf{book-level engineering proposals}. They must not be attributed to BDI, MIDCA, CCA, or Fan \& Lan unless the cited source explicitly provides the mechanism.

\subsubsection{0.4 A distinction that governs the entire chapter}\label{a-distinction-that-governs-the-entire-chapter}

The chapter must keep three authority domains separate:

\begin{itemize}
\tightlist
\item
  \textbf{Goal authority:} who may define or amend the intended outcome and its constraints.
\item
  \textbf{Method authority:} who may continue, delegate, suspend, abandon, or replace the method used to pursue the authorized outcome.
\item
  \textbf{Proposal authority:} what the model or another component may suggest, without thereby authorizing or committing the change.
\end{itemize}

A method can change while the goal remains invariant. A goal change is not a method change and must follow the principal's authorization path.

\subsubsection{0.5 Teaching sequence}\label{teaching-sequence}

Use problem first, concrete intuition, term, precise mechanism, worked example, and boundary. Teach BDI commitment before MIDCA meta-control; then compare modern systems; only after that derive the proposed runtime contracts. Do not start with a list of fields or an implementation diagram without first establishing the problem those fields solve.

\begin{center}\rule{0.5\linewidth}{0.5pt}\end{center}

\subsection{1. Opening: the agent that remembers the task but forgets why it is committed}\label{opening-the-agent-that-remembers-the-task-but-forgets-why-it-is-committed}

\subsubsection{1.1 {[}Book example{]} Product X payment-status incident}\label{book-example-product-x-payment-status-incident}

A user asks an agent to determine whether a payment has settled and report the verified status. A prior action times out. The model proposes a retry, but the external processor may already have accepted the first request. The agent must continue pursuing the same user-authorized goal while reconsidering the unsafe method.

The opening will make the constraints concrete:

\begin{itemize}
\tightlist
\item
  the goal is to establish and communicate the verified payment status;
\item
  the agent must not create a duplicate payment;
\item
  no refund, cancellation, or new payment may be initiated without its own authorization;
\item
  a timeout is evidence about an execution attempt, not proof that the goal is impossible or achieved;
\item
  the runtime may switch from retry-oriented execution to read-only reconciliation without changing the goal;
\item
  if the evidence remains ambiguous, the intention may be suspended or escalated rather than silently abandoned.
\end{itemize}

\subsubsection{1.2 The failure of a task label}\label{the-failure-of-a-task-label}

A task label or prompt fragment may preserve wording but does not necessarily preserve a commitment's authority, rationale, continuation conditions, or termination rules across process restarts and model calls. Repeating the user request in the prompt is not equivalent to a persistent intention record.

\subsubsection{1.3 The failure of reflection alone}\label{the-failure-of-reflection-alone}

A model can notice that a plan is failing and propose another plan. Unless a runtime component owns the committed goal and has authority to accept or reject the proposal, reflection can remain commentary. A persuasive self-critique is not itself a control transition.

\subsubsection{1.4 Learning outcomes}\label{learning-outcomes}

By the end of the chapter, the reader should be able to:

\begin{itemize}
\tightlist
\item
  explain what BDI intentions add beyond beliefs, desires, plans, and task text;
\item
  distinguish commitment termination policies from the timing of reconsideration;
\item
  reconstruct MIDCA's object-level and meta-level relation and explain why diagnosis must have a consequential control path;
\item
  describe the B4 hypothesis without inflating its scope;
\item
  distinguish a new plan, a new method, a delegated method, a suspended intention, and an amended goal;
\item
  design a persistent intention subsystem with typed contracts and explicit authority boundaries;
\item
  implement a method switch while preserving the authorized goal;
\item
  evaluate goal drift, stale commitment, oscillation, premature stopping, and the causal value of method changes.
\end{itemize}

\begin{center}\rule{0.5\linewidth}{0.5pt}\end{center}

\subsection{2. Locate B4 in the Fan \& Lan review}\label{locate-b4-in-the-fan-lan-review}

\subsubsection{2.1 Runtime insertion point}\label{runtime-insertion-point}

Fan \& Lan locate B4 at the \textbf{workflow/governor boundary}. The proposed control law is that an intention persists until explicit reconsideration, while diagnosis may switch methods. Explain why this is neither merely a planner feature nor merely an execution monitor.

\subsubsection{2.2 What the persistent intention record contains}\label{what-the-persistent-intention-record-contains}

Teach the source's stated fields without silently replacing them:

\begin{itemize}
\tightlist
\item
  committed goal;
\item
  rationale;
\item
  plan;
\item
  continuation condition;
\item
  reconsideration triggers;
\item
  authority;
\item
  termination status.
\end{itemize}

Explain the meaning of each field and the type of question it answers. Do not imply that Fan \& Lan specify a database schema or exact serialization format.

\subsubsection{2.3 The governor's decision set}\label{the-governors-decision-set}

Teach the five actions explicitly named in the B4 description:

\begin{itemize}
\tightlist
\item
  continue;
\item
  delegate;
\item
  suspend;
\item
  abandon;
\item
  change method.
\end{itemize}

Relate these to intention lifecycle semantics, but do not make them look like an exhaustive universal policy language. Additional terminal statuses such as achieved or impossible may be used in the book's implementation proposal where their provenance is explicit.

\subsubsection{2.4 Why B4 is not simply BDI relabelled}\label{why-b4-is-not-simply-bdi-relabelled}

BDI supplies the historical commitment structure. B4 asks for a portable runtime lifecycle plus explicit live method authority in contemporary agents. It therefore combines persistent normative state with an operative reconsideration and control edge.

\subsubsection{2.5 Why B4 is not simply reflection}\label{why-b4-is-not-simply-reflection}

Reflection can detect a problem, assess a plan, or propose a replacement. B4 additionally needs an authoritative object representing what remains committed and a governor that can make a permitted state transition.

\subsubsection{2.6 Why B4 is not just CCA or an intent graph}\label{why-b4-is-not-just-cca-or-an-intent-graph}

CCA is a security-oriented lifecycle supervision baseline in Fan \& Lan's review. It contributes intent-graph and deviation-adjudication ideas, but the review identifies its domain-specific scope and the missing general intention lifecycle actions. Do not claim that an intent graph alone implements the complete B4 commitment policy.

\subsubsection{2.7 B4 dependencies and boundaries}\label{b4-dependencies-and-boundaries}

\begin{itemize}
\tightlist
\item
  B1 can supply memories and evidence relevant to reconsideration; it does not own the intention or decide that it should be abandoned.
\item
  B2 can recover from a typed failure within a bounded substate; B4 supplies the parent goal and commitment boundary that recovery must preserve.
\item
  B3 can broadcast evidence or monitor signals that trigger reconsideration; workspace admission does not itself grant method authority.
\item
  B5/GraSP can provide a verified skill or fallback method; B4 determines whether the method remains authorized for the committed goal.
\item
  B6 controls uncertainty, resource allocation, interruption, tool authorization, and stopping; B4 controls persistence and reconsideration of the intention. Avoid teaching B6 prematurely.
\end{itemize}

\subsubsection{2.8 Evidence and migration depth remain separate}\label{evidence-and-migration-depth-remain-separate}

Explain E1-E4 and D0-D4 only as defined in Chapter 11. In this chapter, report the B4 case labels as two separate coordinates, not a composite score:

\begin{itemize}
\tightlist
\item
  \textbf{R5a: hybrid and self-aware BDI-LLM - E1/D3.} Beliefs, goals, intentions, plans, and explicit BDI control are preserved; a portable cross-domain reconsideration lifecycle remains missing.
\item
  \textbf{R5b: Devil's Advocate - E3/D3.} Plan/subtask state, reflection, backtracking, and strategy revision are preserved; a durable authoritative intention and explicit continuation/suspension/abandonment rules remain missing.
\item
  \textbf{R5c: Cognitive Control Architecture - E2/D3.} Intent graph, trajectory deviation, and tiered adjudication are preserved; general-purpose lifecycle authority beyond the security-specific setting remains missing.
\item
  \textbf{R6: ADAS, AFlow, and MetaReflection - E3/D4.} Candidate or learned methods/workflows are informed by design-time or cross-trial experience; live in-run diagnosis and method-switch authority over a persistent intention remain missing.
\end{itemize}

The labels reproduce Fan \& Lan's reviewed records. They are not a judgement of overall system quality, and a higher D value does not mean that the system has solved B4.\FanLanCite{Appendix F, R5a--R5c and R6; Table VI}

\begin{center}\rule{0.5\linewidth}{0.5pt}\end{center}

\subsection{3. Historical foundation I: what BDI means by intention and commitment}\label{historical-foundation-i-what-bdi-means-by-intention-and-commitment}

\subsubsection{3.1 Beliefs, desires, and intentions as distinct attitudes}\label{beliefs-desires-and-intentions-as-distinct-attitudes}

Explain the distinction in the historical BDI account:

\begin{itemize}
\tightlist
\item
  beliefs represent what the agent takes to be the case;
\item
  desires represent states of affairs the agent is motivated to bring about;
\item
  intentions represent selected commitments that constrain subsequent practical reasoning.
\end{itemize}

Do not teach intentions as synonyms for a to-do list, arbitrary plan text, or all goals the user has expressed. A BDI intention has a role in the agent's ongoing deliberation and action selection.

\subsubsection{3.2 The practical reasoning loop}\label{the-practical-reasoning-loop}

Reconstruct the simplified BDI sequence from the cited BDI work:

\begin{enumerate}
\def\labelenumi{\arabic{enumi}.}
\tightlist
\item
  update beliefs from observations and events;
\item
  generate options in relation to desires;
\item
  deliberate and filter options into intentions;
\item
  execute plans associated with intentions;
\item
  update beliefs and intention state after new observations;
\item
  release intentions when their applicable commitment conditions are met.
\end{enumerate}

The chapter must state that this is a teaching reconstruction of the practical loop, not a universal runtime API or a requirement that every modern agent implement the same internal data structures.

\subsubsection{3.3 Why intention matters under bounded deliberation}\label{why-intention-matters-under-bounded-deliberation}

Commitment reduces the need to reconsider every alternative at every step. It supports coherent multi-step behavior, but it can also keep an agent pursuing an obsolete or impossible plan. The mechanism therefore manages a tension between stability and responsiveness.

\subsubsection{3.4 Ends and means must not collapse}\label{ends-and-means-must-not-collapse}

An intention is associated with an intended end and may be pursued through plans or means. The chapter should distinguish preserving the authorized outcome from preserving one specific plan. A change in means does not necessarily imply a change in end; silently changing the end is a different and higher-authority operation.

\subsubsection{3.5 Historical boundary}\label{historical-boundary}

The chapter should explain the BDI contribution as Fan \& Lan reconstruct it and use Rao and Georgeff as a historical cross-check. It should not expand into a full survey of philosophical theories of intention or introduce additional theoretical frameworks outside the primary paper's cited examples.

\begin{center}\rule{0.5\linewidth}{0.5pt}\end{center}

\subsection{4. Historical foundation II: commitment policies and reconsideration schedules are different axes}\label{historical-foundation-ii-commitment-policies-and-reconsideration-schedules-are-different-axes}

\subsubsection{4.1 Why the distinction matters}\label{why-the-distinction-matters}

A common teaching error is to list blind, single-minded, open-minded, bold, and cautious as if they were five members of one interchangeable taxonomy. The chapter must separate \textbf{when an intention may be terminated} from \textbf{how often the agent reconsiders its current options}.

\subsubsection{4.2 Commitment termination policies}\label{commitment-termination-policies}

Explain the three familiar BDI commitment policies with their conditions:

\begin{itemize}
\tightlist
\item
  \textbf{Blind commitment:} retain the intention until the agent believes the intended state has been achieved; conflicting changes in beliefs or desires do not ordinarily release the commitment under this policy.
\item
  \textbf{Single-minded commitment:} retain the intention until the goal is achieved or the agent no longer believes it can be achieved; belief changes can therefore make the commitment terminable.
\item
  \textbf{Open-minded commitment:} retain the intention until achievement or until the relevant belief/desire conditions no longer support continuing to pursue it; the exact condition should follow the formalization being explained.
\end{itemize}

Use Rao and Georgeff's formal/practical explanation to qualify the exact termination conditions. Avoid treating the labels as natural-language personality traits.\autocite{rao1991modelingbdi,rao1995bdi}

\subsubsection{4.3 Reconsideration schedules: bold and cautious}\label{reconsideration-schedules-bold-and-cautious}

Kinny and Georgeff use bold and cautious to describe how frequently options are reconsidered during execution:

\begin{itemize}
\tightlist
\item
  a bold agent does not reconsider options until the current plan has been executed in its entirety;
\item
  a cautious agent reconsiders after each new option/step;
\item
  intermediate schedules can limit the number of plan steps before reconsideration.
\end{itemize}

These describe a schedule or frequency of reconsideration, not the same thing as the three commitment termination conditions in Section 4.2.\autocite{kinny1991commitment}

\subsubsection{4.4 Why neither extreme is universally best}\label{why-neither-extreme-is-universally-best}

Use Kinny and Georgeff's Tileworld experiment to show the cost/benefit trade-off. Frequent replanning consumes deliberation time; infrequent replanning can pursue vanished targets or miss useful changes. The right behavior depends on environmental dynamism, planning cost, the event, and the plan being executed.

\subsubsection{4.5 Event-sensitive reconsideration}\label{event-sensitive-reconsideration}

The historical results support selective response to relevant environmental events rather than reacting to every change. Teach the distinction among:

\begin{itemize}
\tightlist
\item
  a change that makes the goal impossible;
\item
  a change that invalidates the current plan or method;
\item
  a change that merely presents a tempting alternative;
\item
  a change that is irrelevant to the current commitment.
\end{itemize}

The chapter should not invent a universal reconsideration threshold. It should explain why relevance and deliberation cost belong in the policy decision.

\subsubsection{4.6 What B4 adds to the classical distinction}\label{what-b4-adds-to-the-classical-distinction}

B4 must persist the commitment record and expose explicit, inspectable reconsideration triggers and authority in a modern runtime. It also adds live method switching as a consequential control action, not just an internal decision to think again.

\begin{center}\rule{0.5\linewidth}{0.5pt}\end{center}

\subsection{5. Historical foundation III: MIDCA and authority over the method}\label{historical-foundation-iii-midca-and-authority-over-the-method}

\subsubsection{5.1 Object-level and meta-level cycles}\label{object-level-and-meta-level-cycles}

Explain MIDCA's dual-cycle organization:

\begin{itemize}
\tightlist
\item
  the object-level cycle performs ordinary task reasoning and action;
\item
  the meta-level cycle observes traces of object-level activity and reasons about how the object-level process is working.
\end{itemize}

Use the source's note-assess-guide structure:

\begin{enumerate}
\def\labelenumi{\arabic{enumi}.}
\tightlist
\item
  note detects a discrepancy or unexpected condition;
\item
  assess hypothesizes a cause;
\item
  guide selects a response.
\end{enumerate}

\subsubsection{5.2 What the meta-controller may change}\label{what-the-meta-controller-may-change}

Fan \& Lan summarize MIDCA's meta-controller as potentially changing the planner, strategy, goal handling, or another object-level method. The important architectural property is not that a system logs its reasoning; it is that diagnosis can select a different method during execution.

\subsubsection{5.3 Method authority is not goal authority}\label{method-authority-is-not-goal-authority}

Teach the key distinction:

\begin{itemize}
\tightlist
\item
  changing the method changes how an authorized goal will be pursued;
\item
  changing the goal changes what outcome the agent is committed to;
\item
  the latter must not be smuggled into the former.
\end{itemize}

The B4 implementation proposal will give the method governor a bounded authority to select an approved alternative, while goal amendment remains with the principal or another explicitly authorized actor.

\subsubsection{5.4 Monitoring without authority is not meta-control}\label{monitoring-without-authority-is-not-meta-control}

A dashboard, trace, evaluator, reflection prompt, or guardrail can identify a problem. It does not implement method authority unless its output can trigger a governed state transition. Conversely, an unconstrained model rewrite is not safe meta-control merely because it changes a plan.

\subsubsection{5.5 Historical boundary}\label{historical-boundary-1}

Do not collapse BDI and MIDCA into one historical architecture. BDI explains intention commitment; MIDCA contributes a meta-level model of diagnosing and changing methods. B4's proposed combination is an analytical synthesis in Fan \& Lan, not a claim that BDI alone already provides MIDCA's meta-controller.\autocite{cox2016midca}\FanLanCite{\S III-C; Figure 3; Appendix A, MIDCA}

\begin{center}\rule{0.5\linewidth}{0.5pt}\end{center}

\subsection{6. Modern evidence: what contemporary systems preserve and what they omit}\label{modern-evidence-what-contemporary-systems-preserve-and-what-they-omit}

\subsubsection{6.1 Hybrid BDI-LLM systems (R5a)}\label{hybrid-bdi-llm-systems-r5a}

Use the child-helpline training agent as a domain-specific example. Explain the BDI control framework and the LLM's role in intent recognition, response generation, and bypass behavior as reported in its source. State Fan \& Lan's assessment: explicit BDI control and persistent structured commitment are preserved, but a portable cross-domain reconsideration lifecycle remains missing. Label E1/D3 as the paper's record.\autocite{owayyed2025controlled,lawton2026selfawarebdi}\FanLanCite{Appendix F, R5a}

\subsubsection{6.2 Self-aware BDI-LLM (R5a)}\label{self-aware-bdi-llm-r5a}

Describe the 2026 University of Utah technical report at the level relevant to B4: a BDI deliberation loop, machine-consumable schemas, plan matching, verification tiers, parameterized plans, and a preliminary evaluation. Clearly identify it as a technical report with preliminary results rather than a settled production standard. Explain what it contributes to structured intention handling and what still needs a general method-governance lifecycle.

\subsubsection{6.3 Devil's Advocate (R5b)}\label{devils-advocate-r5b}

Explain the three interventions in the cited paper:

\begin{itemize}
\tightlist
\item
  anticipatory reflection before action, considering failure and alternatives;
\item
  post-action alignment with subtask objectives and backtracking/remedy;
\item
  overall plan review after completion.
\end{itemize}

This is a meaningful plan-revision baseline. It must not be misrepresented as a durable, authoritative intention registry with explicit continue/delegate/suspend/abandon policy. Preserve Fan \& Lan's E3/D3 record.\autocite{wang2024devilsadvocate}\FanLanCite{Appendix F, R5b}

\subsubsection{6.4 Cognitive Control Architecture (R5c)}\label{cognitive-control-architecture-r5c}

Explain the intent graph and tiered adjudicator at a high level, only to the extent needed for B4 comparison. It is security-oriented: a deviation from an expected action trajectory triggers additional adjudication. Preserve Fan \& Lan's E2/D3 record and its boundary: security-specific trajectory integrity does not automatically supply a domain-general intention lifecycle.\autocite{liang2025cca}\FanLanCite{Appendix F, R5c}

\subsubsection{6.5 Goal-drift evidence}\label{goal-drift-evidence}

Use the goal-drift study as a diagnostic warning that initially specified objectives can drift during long-horizon interactions under competing pressures. Do not present a diagnostic study as a persistence controller or claim that a B4 governor is already validated by it.\autocite{arike2025goaldrift}

\subsubsection{6.6 Premature-commitment evidence}\label{premature-commitment-evidence}

Use the 2026 premature-commitment study to distinguish settling on a trajectory from being correct. Its signal is diagnostic; the reported work explicitly does not establish that hidden-state convergence is a correctness signal. Do not turn this metric into the proposed B4 governor's sole trigger.\autocite{mehta2026premature}

\subsubsection{6.7 ADAS, AFlow, and MetaReflection (R6)}\label{adas-aflow-and-metareflection-r6}

Explain their relation to redesigning or improving agent/workflow methods across design time or trials. Preserve Fan \& Lan's E3/D4 classification and the boundary: learning better methods across trials is not the same as an in-run governor having authority to change the method for a persistent intention while the goal remains fixed.\autocite{hu2024adas,zhang2024aflow,gupta2024metareflection}\FanLanCite{Appendix F, R6}

\subsubsection{6.8 Comparative table}\label{comparative-table}

Compare each system by its actual state object, trigger, control action, persistence boundary, what it contributes to B4, and the missing lifecycle edge. Keep E and D as separate columns. Do not infer unspecified internals from a system's name or popularity.

\subsubsection{6.9 What the comparison establishes - and what it does not}\label{what-the-comparison-establishes---and-what-it-does-not}

The reviewed systems demonstrate important pieces: explicit BDI commitments, reflection-triggered revision, intent-graph adjudication, goal-drift diagnostics, premature-commitment diagnostics, and cross-trial method improvement. They do not, as a set, establish the full portable B4 coupling. The chapter should avoid claiming that no existing system has any relevant feature or that B4 is validated merely because its components have precedents.\FanLanCite{Appendix F, R5a--R5c and R6}

\begin{center}\rule{0.5\linewidth}{0.5pt}\end{center}

\subsection{7. B4 reconstructed through the seven mechanism fields}\label{b4-reconstructed-through-the-seven-mechanism-fields}

Use Chapter 10's field names and order exactly. Do not add an eighth field or reinterpret the E/D labels as part of the tuple.

\subsubsection{7.1 State (S)}\label{state-s}

Persistent intention, committed goal, rationale, active plan/method, continuation conditions, reconsideration triggers, authority, termination status, beliefs/evidence relevant to goal achievement, and method execution trace. Mark details not specified by Fan \& Lan as unspecified in the source and reserve schema elaboration for Section 9.

\subsubsection{7.2 Control (C)}\label{control-c}

Commitment policy, deliberation/option filtering, reconsideration governor, MIDCA-like meta-controller, and the authority boundary between proposal and commitment. Identify which actor is documented in the historical sources and which proposed runtime component is book-designed.

\subsubsection{7.3 Transition (T)}\label{transition-t}

Observation, goal success, perceived impossibility, material belief change, plan/method failure, goal relevance change, discrepancy diagnosis, or an explicit authority event can trigger reconsideration. Do not imply every event must cause a method change.

\subsubsection{7.4 Persistence (P)}\label{persistence-p}

The intention record survives individual model calls, planning steps, tool invocations, and process restarts as long as its commitment conditions remain valid. Method assignment may change while the authorized goal identity and constraints remain stable.

\subsubsection{7.5 Failure (F)}\label{failure-f}

Stale commitment, goal drift, method thrashing/oscillation, premature stopping, false impossibility, plan/method change without diagnostic value, authority confusion, and silent goal substitution.

\subsubsection{7.6 Learning (L)}\label{learning-l}

Separate runtime reconsideration from across-trial policy or method learning. Explain which reviewed systems update future workflow designs or instructions, and why that is not automatically live B4 authority. Any proposed learned reconsideration policy must be audited, versioned, and gated.

\subsubsection{7.7 Resources and governance (R)}\label{resources-and-governance-r}

Deliberation cost, replanning cost, time, resource budgets, delegated scope, authorization, auditability, and safe interruption. Defer the full cross-stage uncertainty and resource governor to Chapter 16 (B6).

\begin{center}\rule{0.5\linewidth}{0.5pt}\end{center}

\subsection{8. {[}Book example{]} Product X: preserve the payment-status goal, change the method}\label{book-example-product-x-preserve-the-payment-status-goal-change-the-method}

\subsubsection{8.1 Authorized goal and invariants}\label{authorized-goal-and-invariants}

The user asks for verified payment status. The immutable goal specification states the outcome, acceptable evidence, no-duplicate-payment constraint, and the fact that a refund or new payment requires separate authorization.

\subsubsection{8.2 Initial intention and method}\label{initial-intention-and-method}

Create an intention to reconcile the ledger and payment processor. The initial method may inspect the local ledger, query the processor status, and compare transaction identifiers. The intention record stores why the goal remains active and what events warrant review.

\subsubsection{8.3 A timeout is not a terminal result}\label{a-timeout-is-not-a-terminal-result}

A request timeout does not establish whether a side effect occurred. The method may have failed; the goal has not necessarily become impossible or satisfied. The model's suggestion to retry is only a proposal.

\subsubsection{8.4 Trigger and diagnosis}\label{trigger-and-diagnosis}

The execution monitor records an unknown external result. The reconsideration detector identifies a material method failure and evidence ambiguity. The diagnosis component determines which facts are known, unknown, or stale and whether the current method can safely proceed.

\subsubsection{8.5 Governor decision: switch method, not goal}\label{governor-decision-switch-method-not-goal}

The governor rejects an unguarded retry, preserves the goal specification, and switches to read-only reconciliation using a processor idempotency lookup or settlement query. If authoritative status cannot be obtained, the intention is suspended or escalated. The governor records why this method is permitted and what evidence would permit resumption.

\subsubsection{8.6 Delegation without goal transfer}\label{delegation-without-goal-transfer}

A bounded read-only reconciliation task may be delegated to a payment-operations agent or human. The delegation contract carries parent intention ID, scope, permitted observations, prohibited actions, expiry, required evidence, and a return schema. The delegate does not acquire permission to amend the user's goal or initiate an external financial operation.

\subsubsection{8.7 Completion and audit}\label{completion-and-audit}

When authoritative evidence resolves the status, the runtime records whether the goal is satisfied and how. The history makes the original commitment, trigger, diagnosis, method switch, delegation, evidence, and termination decision reconstructable.

\subsubsection{8.8 Counterexamples}\label{counterexamples}

Show at least four failure variants: blind retry creates duplicate payment; agent changes the goal to ``attempt another payment''; reflection proposes a switch but nothing authorizes it; or governor switches repeatedly between methods without new evidence.

\begin{center}\rule{0.5\linewidth}{0.5pt}\end{center}

\subsection{9. Theory-to-practice bridge: a reference intention and method-governance subsystem}\label{theory-to-practice-bridge-a-reference-intention-and-method-governance-subsystem}

This section is mandatory and must be one of the chapter's major sections. It translates the B4 control relation into a proposed implementation without claiming that the source paper provides these software interfaces.

\subsubsection{9.1 Reference architecture}\label{reference-architecture}

Describe the roles and their trust boundaries:

\begin{enumerate}
\def\labelenumi{\arabic{enumi}.}
\tightlist
\item
  Principal/Goal API: accepts and authenticates the user-authorized goal and later goal amendments.
\item
  Goal specification store: persists the immutable, versioned goal and its constraints.
\item
  Intention registry: persists active intention records and lifecycle state.
\item
  Belief/evidence store: records observations, provenance, timestamps, and confidence/verification state.
\item
  Execution monitor: observes plan and tool outcomes without deciding policy by itself.
\item
  Reconsideration detector: converts configured triggers into structured events.
\item
  Diagnosis service: examines evidence and proposes a cause/impact assessment; model reasoning is advisory.
\item
  Method registry: lists permitted methods with version, capabilities, preconditions, postconditions, cost, risk class, and supported goal types.
\item
  Method-authority governor: selects among allowed lifecycle actions under the intention's authority policy.
\item
  Authorization and safety gate: checks principal consent, tool permissions, external effects, and prohibited transitions independently.
\item
  Executor/delegation adapter: invokes the current method or scoped delegate.
\item
  Append-only event and decision ledger: records proposals, diagnoses, approvals, state transitions, outcomes, and version references.
\end{enumerate}

\subsubsection{9.2 Typed data contracts}\label{typed-data-contracts}

The chapter should specify fields and invariants for each contract:

\begin{itemize}
\tightlist
\item
  \texttt{GoalSpec}: goal ID, principal/owner, authorized outcome, immutable constraints, acceptance criteria, permitted side effects, version, content hash, and authorization evidence.
\item
  \texttt{IntentionRecord}: intention ID, goal reference/hash, rationale and evidence references, plan reference, active method assignment, continuation condition, reconsideration triggers, authority policy reference, termination status, timestamps, and version.
\item
  \texttt{MethodDescriptor}: method ID/version, eligible goal types, capabilities, prerequisites, preconditions, postconditions, expected side effects, cost/risk classification, and required authority.
\item
  \texttt{MethodAssignment}: intention ID, method descriptor/version, invocation ID, plan state, assignment reason, validity interval, and status.
\item
  \texttt{ReconsiderationEvent}: event ID, intention/version, trigger type, evidence IDs, provenance, observation time, goal-relevance result, and diagnosis reference.
\item
  \texttt{GovernorDecision}: decision ID, action, prior and next lifecycle state, prior and next method, rationale/evidence, authority path, goal hash, expected state version, and timestamp.
\item
  \texttt{DelegationContract}: parent intention ID, delegated objective within the parent goal, scope, allowed/prohibited capabilities, evidence obligations, expiry, cancellation rule, and return contract.
\item
  \texttt{AuthorityPolicy}: principal, governor permissions, model proposal-only permissions, human approval requirements, safety vetoes, and audit requirements.
\end{itemize}

Explain which fields are source-derived abstractions from B4 and which are implementation-specific additions. Do not store model-generated rationale as if it were the principal's authorization evidence.

\subsubsection{9.3 Separate goal, method, and proposal authorities}\label{separate-goal-method-and-proposal-authorities}

Define the proposed permission matrix:

\begin{itemize}
\tightlist
\item
  \textbf{Principal/user:} may define or amend the goal through authenticated authorization; may pre-authorize a bounded set of method choices; external side effects remain subject to explicit consent and action policy.
\item
  \textbf{Language model:} may propose methods or diagnoses; cannot commit a proposal, amend the goal, or authorize an external side effect.
\item
  \textbf{Reconsideration governor:} may select only from the intention's permitted method set and may request alternatives; its method authority does not include goal amendment or external-action approval.
\item
  \textbf{Safety/authorization gate:} may block, require review, or enforce the separate policy for external side effects; it does not invent a new goal.
\item
  \textbf{Executor/delegate:} may execute only the assigned method and scoped capabilities; it may return evidence and proposals but cannot enlarge its authority.
\end{itemize}

The exact matrix is a book-level design proposal. It should be related to the B4 requirement that reconsideration authority be explicit, while noting that external-action authorization is also a B6 concern.

\subsubsection{9.4 Intention lifecycle state machine}\label{intention-lifecycle-state-machine}

Propose a versioned state machine with:

\begin{itemize}
\tightlist
\item
  \texttt{PROPOSED}: candidate intention exists but is not committed;
\item
  \texttt{AUTHORIZED}: principal/policy has accepted the goal and allowed scope;
\item
  \texttt{ACTIVE}: the intention is being pursued through a current method;
\item
  \texttt{RECONSIDERATION\_PENDING}: a trigger requires diagnosis and a governor decision;
\item
  \texttt{METHOD\_SWITCH\_PENDING}: a replacement method is being checked and committed;
\item
  \texttt{DELEGATED}: a child assignment is active under a contract;
\item
  \texttt{SUSPENDED}: no method may proceed until a defined resume condition/authority event;
\item
  \texttt{SATISFIED}, \texttt{IMPOSSIBLE}, and \texttt{ABANDONED}: terminal outcomes, with evidence and authority requirements.
\end{itemize}

Define legal transitions and guards. A method switch changes the method assignment but preserves the \texttt{GoalSpec} ID/hash and immutable goal constraints. Goal amendment is a separate principal-authorized operation that creates a new goal version and a recorded relationship to the previous intention. A terminal intention is not silently reopened by changing a prompt; a new or explicitly reauthorized intention is required.

\subsubsection{9.5 Reconsideration trigger taxonomy}\label{reconsideration-trigger-taxonomy}

Cover success evidence, impossibility evidence, material belief/world change, goal-relevance change, method precondition failure, repeated lack of progress, unexpected side-effect evidence, delegate return, and explicit principal instruction. Separate hard invalidators from soft reconsideration signals. A low-value alternative appearing is not sufficient by itself to abandon a commitment.

\subsubsection{9.6 Diagnosis and decision semantics}\label{diagnosis-and-decision-semantics}

A diagnosis should report observed evidence, provenance, current method assumptions, what changed, which conditions remain valid, goal feasibility, goal relevance, and uncertainty. The governor then chooses a permitted action. The LLM may generate candidate explanations or methods, but it cannot label an intention satisfied/impossible or amend the goal without the designated authority and evidence checks.

\subsubsection{9.7 Distinguish plan revision from method switch}\label{distinguish-plan-revision-from-method-switch}

A plan may be revised within the same method family while preserving the method authority contract. A method switch changes the mechanism or strategy used to pursue the goal and should be logged as such. The chapter must define the distinction operationally and explain why otherwise every minor plan edit could look like a B4 transition.

\subsubsection{9.8 Persistence, transactions, and concurrency}\label{persistence-transactions-and-concurrency}

Propose:

\begin{itemize}
\tightlist
\item
  append-only intention and reconsideration events;
\item
  versioned snapshots for fast recovery;
\item
  compare-and-swap or optimistic concurrency on intention versions;
\item
  event IDs and idempotent processing for repeated triggers;
\item
  transactional updates for method assignment and state transition;
\item
  outbox/inbox patterns or equivalent durable delivery for asynchronous method/delegate calls;
\item
  revalidation of evidence and method preconditions before commit;
\item
  restart recovery that reconstructs the active commitment and pending work;
\item
  conflict handling when two monitor events race to reconsider the same intention.
\end{itemize}

Explain the invariant: a stale governor decision cannot overwrite a newer intention version or change its goal hash.

\subsubsection{9.9 Method switch and external side effects}\label{method-switch-and-external-side-effects}

Changing a method does not reverse actions already performed. A payment submission, message, file write, or external API mutation must be reconciled or compensated through its own authorized procedure. Do not describe method switch as rollback of real-world effects. Require idempotency and operation IDs where supported, record uncertain results, and route compensation through the external-action authorization policy.

\subsubsection{9.10 Delegation and suspension}\label{delegation-and-suspension}

A delegate receives bounded work, not the entire parent authority. Define scope, tool allowlist, deadlines, expected evidence, return format, cancellation behavior, escalation path, and how the parent intention resumes. Suspension must specify what is paused, what events can resume it, and whether a human decision is required.

\subsubsection{9.11 Anti-oscillation and commitment safeguards}\label{anti-oscillation-and-commitment-safeguards}

Propose evidence thresholds, hysteresis, cooldowns, switch budgets, minimum expected diagnostic value, and a record of rejected alternatives. These safeguards must not prevent immediate suspension for hard safety violations or falsely keep an impossible method alive. Distinguish a stable goal from a frozen plan.

\subsubsection{9.12 Java-like interfaces and orchestration pseudocode}\label{java-like-interfaces-and-orchestration-pseudocode}

Provide implementation-oriented but non-framework-specific Java records/interfaces. The chapter should show, at minimum:

\begin{itemize}
\tightlist
\item
  immutable \texttt{GoalSpec} and versioned \texttt{IntentionRecord};
\item
  \texttt{ReconsiderationEvent} and typed \texttt{GovernorDecision};
\item
  \texttt{MethodDescriptor} and \texttt{DelegationContract};
\item
  \texttt{IntentionRepository} with compare-and-swap semantics;
\item
  \texttt{MethodRegistry} and \texttt{MethodExecutor} interfaces;
\item
  a \texttt{ReconsiderationGovernor} that returns a typed decision rather than mutating state implicitly.
\end{itemize}

Pseudocode should show: load current version, validate trigger, gather evidence, diagnose, ask/obtain proposal, validate against authority, commit the permitted transition atomically, dispatch the method/delegate, and append outcomes. Keep the goal hash invariant explicit.

\subsubsection{9.13 Incremental implementation sequence}\label{incremental-implementation-sequence}

\begin{enumerate}
\def\labelenumi{\arabic{enumi}.}
\tightlist
\item
  Persist the principal-authorized goal separately from prompt text.
\item
  Add a minimal intention record and lifecycle with explicit completion/impossibility rules.
\item
  Instrument execution and reconsideration events.
\item
  Add an allowlisted method registry and versioned assignments.
\item
  Implement the governor and separate authorization gate.
\item
  Add suspension and bounded delegation.
\item
  Add append-only audit, restart recovery, concurrency and idempotency.
\item
  Introduce anti-oscillation policy and learned method recommendations only after the non-learning control path is measurable.
\end{enumerate}

Each step should have a demonstrable acceptance test. Do not begin with an unconstrained self-modifying planner.

\subsubsection{9.14 Testable implementation scenarios}\label{testable-implementation-scenarios}

Include tests for: goal hash unchanged after method switch; model cannot amend goal; principal amendment produces new goal version; satisfied intention terminates; method precondition invalidation triggers review; stale diagnosis loses compare-and-swap; duplicate trigger is idempotent; delegate exceeds scope and is blocked; suspension halts the method; a delegate result resumes only after validation; method switch does not erase prior external side effects; repeated low-value alternatives do not cause oscillation; impossible goal is not kept active indefinitely.

\begin{center}\rule{0.5\linewidth}{0.5pt}\end{center}

\subsection{10. Failure modes, safety boundaries, and operational guidance}\label{failure-modes-safety-boundaries-and-operational-guidance}

\subsubsection{10.1 Stale commitment}\label{stale-commitment}

The agent continues a method even after key assumptions or the goal's feasibility have changed. Show how a continuation condition and event-sensitive reconsideration address it.

\subsubsection{10.2 Goal drift and silent substitution}\label{goal-drift-and-silent-substitution}

The runtime's declared goal remains one thing while execution increasingly optimizes another. Require goal identity/hash, acceptance criteria, provenance, and principal-authorized amendment records.

\subsubsection{10.3 Method oscillation}\label{method-oscillation}

The governor switches repeatedly without new evidence or improved expected outcome. Use event relevance, hysteresis, cooldown, switch budgets, and review metrics.

\subsubsection{10.4 Premature stopping}\label{premature-stopping}

The agent treats uncertainty, one failure, or one negative signal as proof that the goal is impossible. Require typed diagnosis and evidence appropriate to the termination condition.

\subsubsection{10.5 Premature commitment to a wrong interpretation}\label{premature-commitment-to-a-wrong-interpretation}

The agent settles early and then defends a mistaken reading. A stability/convergence signal can be a reason to inspect the trajectory, but the cited premature-commitment diagnostic does not establish correctness; it must not be a sole method-switch trigger.

\subsubsection{10.6 Reflection without control}\label{reflection-without-control}

The system produces a critique but has no authoritative state transition. Add a test proving that the governor decision changes the active method only through a valid state transition.

\subsubsection{10.7 Method authority becomes goal authority}\label{method-authority-becomes-goal-authority}

An LLM or subagent proposes a new goal under the guise of selecting a method. Enforce the immutable goal reference and separate principal authorization path.

\subsubsection{10.8 External side-effect confusion}\label{external-side-effect-confusion}

A method switch, cancellation, or delegation cannot erase or reverse an already executed external action. Preserve operation IDs and require separate compensation authority.

\subsubsection{10.9 Over-monitoring and excessive deliberation}\label{over-monitoring-and-excessive-deliberation}

Reconsidering on every message can consume the resources needed to act. Use event-sensitive triggers and measure replan/deliberation cost.

\subsubsection{10.10 Cross-bundle safety boundaries}\label{cross-bundle-safety-boundaries}

\begin{itemize}
\tightlist
\item
  B2 owns typed recovery substates and resolution return; B4 retains the parent intention and decides whether its method should continue after the resolution.
\item
  B3 may broadcast relevant evidence; B4 owns the commitment decision.
\item
  B6 owns uncertainty propagation, budgets, interruption, tool authorization, and stopping. B4 must not duplicate B6's governor.
\item
  B5 provides eligible skill/fallback methods but does not own the parent goal.
\end{itemize}

\begin{center}\rule{0.5\linewidth}{0.5pt}\end{center}

\subsection{11. Evaluation: does explicit intention and method authority improve long-horizon behavior?}\label{evaluation-does-explicit-intention-and-method-authority-improve-long-horizon-behavior}

\subsubsection{11.1 Preserve Fan \& Lan Appendix D}\label{preserve-fan-lan-appendix-d}

The chapter must reproduce the required B4 trace:

\begin{enumerate}
\def\labelenumi{\arabic{enumi}.}
\tightlist
\item
  intention record;
\item
  reconsideration event;
\item
  diagnosis;
\item
  method switch;
\item
  termination.
\end{enumerate}

The study should hold base model, prompt, tools, environment, and task budget constant across controlled conditions. Preregister the task population, model/dependency versions, experimental unit, stochastic replications, sample size, metrics, non-inferiority/equivalence margin, uncertainty intervals/tests, exclusions, and stopping rule. Do not invent sample-size numbers in advance of a real power analysis.\FanLanCite{Appendix D; Appendix F, B4 record}

\subsubsection{11.2 Primary outcomes}\label{primary-outcomes}

\begin{itemize}
\tightlist
\item
  goal drift / adherence to the originally authorized goal;
\item
  task completion under the stated acceptance criteria;
\item
  correct stopping or continuation decision, including justified impossibility and suspension.
\end{itemize}

\subsubsection{11.3 Safety and efficiency outcomes}\label{safety-and-efficiency-outcomes}

\begin{itemize}
\tightlist
\item
  method-switch and replan frequency;
\item
  oscillation rate and repeated switch-back;
\item
  escalation precision;
\item
  stale-intention rate;
\item
  unauthorized goal-change attempts;
\item
  cost and elapsed time attributable to reconsideration;
\item
  method-switch diagnostic value and post-switch performance.
\end{itemize}

Treat any additional metric definitions as book-proposed operationalizations, not as values already reported by Fan \& Lan.

\subsubsection{11.4 Baselines}\label{baselines}

Compare against: a prompt-only agent with no persistent intention record; a static plan that never reconsiders; a BDI-like commitment baseline without live meta-method authority; reflection-triggered plan revision; and the relevant reviewed modern examples (hybrid/self-aware BDI-LLM, Devil's Advocate, CCA). Use ADAS/AFlow/MetaReflection only as design-time or cross-trial learning comparators, not as direct live-governor equivalents.

\subsubsection{11.5 Required ablations}\label{required-ablations}

Preserve the three source-specified B4 ablations:

\begin{enumerate}
\def\labelenumi{\arabic{enumi}.}
\tightlist
\item
  remove the explicit persistent intention state object;
\item
  remove the authoritative control edge while retaining monitoring/logging;
\item
  remove the learning update while retaining the runtime transition.
\end{enumerate}

Explain what each ablation isolates. Do not merge the three into one generic ``remove memory'' condition.

\subsubsection{11.6 Controlled task suite}\label{controlled-task-suite}

Propose tasks that vary environmental dynamism, method reliability, plan cost, feasibility changes, goal relevance, distractor pressure, and delegated work. Include long-horizon tasks where the same goal remains valid but the original method becomes invalid, and tasks where the goal itself becomes impossible or is explicitly amended. Use matched seeds and deterministic event injection where possible.

\subsubsection{11.7 Falsification criteria}\label{falsification-criteria}

Reject or narrow B4 if: explicit intention state does not reduce goal drift; method switching does not improve outcomes over the strongest simpler baseline; reconsideration oscillates; completion or correct stopping worsens; method changes lack diagnostic value; or the lifecycle adds cost without measurable benefit. Preserve Fan \& Lan's named failure modes: stale commitment, oscillation, premature stopping, and method changes without diagnostic value.\FanLanCite{Table VI; Appendix D, Table X}

\subsubsection{11.8 Threats to validity}\label{threats-to-validity}

Discuss model and environment confounds, goal-drift measurement ambiguity, evaluator bias, false diagnosis labels, benchmark overfitting, domain-specific authority policies, and the difference between method change and plan editing. Keep results bounded to the tested tasks and versions.

\begin{center}\rule{0.5\linewidth}{0.5pt}\end{center}

\subsection{12. Relationship to the other bundles and chapter handoff}\label{relationship-to-the-other-bundles-and-chapter-handoff}

\subsubsection{12.1 Dependency summary}\label{dependency-summary}

Fan \& Lan identify B4 as the commitment boundary that scopes recovery and helps preserve the parent goal during B2. Workspace evidence may trigger B4 reconsideration. B6 governs risk, resources, tool authorization, interruption, and stopping. This chapter should use those relations to explain compositional control without pre-writing Chapters 16 or 18.\FanLanCite{Figure 8; Appendix D, cross-bundle dependencies}

\subsubsection{12.2 What Chapter 15 does not solve}\label{what-chapter-15-does-not-solve}

B4 does not itself supply a typed recovery substate, workspace admission policy, learned memory selector, skill verification DAG, or calibrated cross-stage uncertainty governor. It decides how a persistent intention remains active and how a method may be reconsidered under authority.

\subsubsection{12.3 Handoff to Chapter 16 (B6)}\label{handoff-to-chapter-16-b6}

Chapter 16 will address how uncertainty and resource budgets influence verification, interruption, tool authorization, and stopping across the execution process. Chapter 15 should flag that boundary and stop before teaching B6.

\begin{center}\rule{0.5\linewidth}{0.5pt}\end{center}

\subsection{13. Recap, review questions, and practical exercises}\label{recap-review-questions-and-practical-exercises}

\subsubsection{13.1 Chapter recap}\label{chapter-recap}

The recap should return to the central distinction: stable commitment to an authorized end does not require rigid commitment to one means; method flexibility is safe only when a durable intention and explicit authority structure govern it.

\subsubsection{13.2 Review questions}\label{review-questions}

\begin{enumerate}
\def\labelenumi{\arabic{enumi}.}
\tightlist
\item
  Why is a prompt containing the user goal not equivalent to a persistent intention record?
\item
  How do blind, single-minded, and open-minded commitment differ?
\item
  Why are bold and cautious reconsideration schedules a separate axis?
\item
  What does MIDCA add beyond object-level planning and an execution log?
\item
  What authority must a governor possess to make reflection consequential?
\item
  Why must goal authority be separated from method authority and model proposal authority?
\item
  What does Fan \& Lan's E1/D3 record for BDI-LLM preserve and omit?
\item
  Why is CCA not a complete general B4 implementation?
\item
  How can a method switch preserve a payment-status goal after an uncertain timeout?
\item
  Which ablation tests whether logging alone explains the improvement?
\end{enumerate}

\subsubsection{13.3 Practical exercises}\label{practical-exercises}

\begin{itemize}
\tightlist
\item
  Design an intention record for a travel-booking goal and identify which fields remain stable when the booking method changes.
\item
  Write a transition table that separates method change, suspension, delegation, impossibility, abandonment, and goal amendment.
\item
  Implement a compare-and-swap method-switch operation and test stale-event handling.
\item
  Create an event sequence where reacting to every new option is worse than selective reconsideration.
\item
  Design a goal-drift experiment with a matched prompt-only baseline and an authoritative intention governor.
\item
  Review a hypothetical model-generated proposal that changes the goal and identify the authority check that must reject it.
\end{itemize}

\begin{center}\rule{0.5\linewidth}{0.5pt}\end{center}

\subsection{14. Planned figures and Mermaid source files}\label{planned-figures-and-mermaid-source-files}

All figures will be authored as Mermaid sources in the chapter-specific source folder and included with \texttt{\textbackslash{}BookDiagram} in the chapter. The following list is a drafting plan, not a claim that PNGs already exist.

\begin{enumerate}
\def\labelenumi{\arabic{enumi}.}
\tightlist
\item
  \texttt{bdi-intention-cycle.mmd} --- beliefs, desires, option filtering, intention, plan execution, observation, and commitment release; historical teaching reconstruction.
\item
  \texttt{commitment-versus-reconsideration.mmd} --- two separate axes: termination conditions (blind/single-minded/open-minded) and reconsideration schedule (bold/cautious/intermediate).
\item
  \texttt{midca-dual-cycle.mmd} --- object-level cognitive cycle and meta-level note-assess-guide control.
\item
  \texttt{b4-authority-architecture.mmd} --- principal, goal store, intention registry, monitor, diagnosis, method governor, model proposals, method registry, safety gate, executor, delegation, and ledger.
\item
  \texttt{b4-intention-lifecycle.mmd} --- proposed intention states and guarded transitions.
\item
  \texttt{productx-method-switch-sequence.mmd} --- payment timeout, event, diagnosis, method switch, read-only reconciliation, possible suspension/delegation, verified completion.
\item
  \texttt{b4-outcome-evaluation.mmd} --- required trace, outcomes, baselines, and the three source-specified ablations.
\end{enumerate}

Every source must include \texttt{accTitle} and \texttt{accDescr}, use the shared theme, have a concise label set, and be referenced exactly once from the corresponding chapter figure placement unless a documented repeat is required.

\begin{center}\rule{0.5\linewidth}{0.5pt}\end{center}

\subsection{15. Planned glossary and index integration}\label{planned-glossary-and-index-integration}

\subsubsection{15.1 Glossary candidates}\label{glossary-candidates}

Before adding entries, check the shared glossary for duplicates and reuse stable terms. Candidate entries:

\begin{itemize}
\tightlist
\item
  belief;
\item
  desire;
\item
  intention;
\item
  intention commitment;
\item
  commitment strategy;
\item
  intention reconsideration;
\item
  reconsideration trigger;
\item
  bold reconsideration;
\item
  cautious reconsideration;
\item
  object-level cycle;
\item
  meta-level cycle;
\item
  meta-control;
\item
  method authority;
\item
  method switch;
\item
  goal drift;
\item
  premature commitment;
\item
  continuation condition;
\item
  delegation contract;
\item
  intention lifecycle.
\end{itemize}

Only add terms actually taught in the final chapter. Do not create synonyms for existing glossary entries.

\subsubsection{15.2 Index candidates}\label{index-candidates}

BDI; belief; desire; intention; commitment; commitment policy; blind commitment; single-minded commitment; open-minded commitment; bold reconsideration; cautious reconsideration; goal drift; method authority; method switch; delegation; suspension; abandonment; MIDCA; meta-control; note-assess-guide; persistent intention; premature commitment; goal authority.

\begin{center}\rule{0.5\linewidth}{0.5pt}\end{center}

\subsection{16. Source-fidelity and attribution plan}\label{source-fidelity-and-attribution-plan}

\subsubsection{16.1 Fan \& Lan (2026)}\label{fan-lan-2026}

Primary source for the chapter's B4 residual, historical-to-modern comparison, E/D case labels, insertion point, closest baselines, failure conditions, cross-bundle dependencies, and Appendix D experiment template. Cite the relevant BDI/MIDCA historical reconstruction, B4 subsection, Table VI, Figure 8, R5a-R5c/R6 records, Table X, and Appendix D directly. Do not imply that the authors report a completed B4 implementation or an experiment validating the bundle.

\subsubsection{16.2 Historical BDI sources}\label{historical-bdi-sources}

Use Rao and Georgeff (1991) for formal intention and mental-attitude framing; Rao and Georgeff (1995) for the practical BDI agent architecture and commitment termination conditions; Kinny and Georgeff (1991) for experimental distinctions among bold, cautious, intermediate reconsideration schedules and event-sensitive reaction. These historical sources are cited by Fan \& Lan. Keep commitment termination and reconsideration frequency separate.

\subsubsection{16.3 Historical MIDCA source}\label{historical-midca-source}

Use Cox et al.~(2016) to explain the dual-cycle architecture and note-assess-guide procedure. The meta-level cycle can monitor and guide object-level method changes. Do not claim that MIDCA is identical to BDI or that the book's proposed intention schema appears in MIDCA.

\subsubsection{16.4 Modern sources}\label{modern-sources}

Use only the specific systems and studies cited or discussed by Fan \& Lan. Attribute each described mechanism to its own source; use Fan \& Lan for the E/D mapping. Record publication status accurately: peer-reviewed proceedings/journal papers versus arXiv preprints and technical reports. In particular, the self-aware BDI work is a 2026 technical report with preliminary evaluation; CCA and the premature-commitment work are arXiv preprints unless a newer published version is verified at drafting time.

\subsubsection{16.5 Claim categories}\label{claim-categories}

Use explicit wording to distinguish:

\begin{itemize}
\tightlist
\item
  \textbf{{[}Fan \& Lan{]}} what the review directly reports or proposes;
\item
  \textbf{historical source reconstruction} what the cited BDI/MIDCA source describes;
\item
  \textbf{modern system report} what the named modern paper actually implements or evaluates;
\item
  \textbf{{[}Book example{]}} Product X and any other illustrative scenario;
\item
  \textbf{{[}Engineering decision{]}} proposed records, contracts, state machine, transaction semantics, authority matrix, safeguards, and test plan.
\end{itemize}

Only the book's two formal paper labels \texttt{{[}CoALA{]}} and \texttt{{[}Fan\ \&\ Lan{]}} may be used. Chapter 15 is Part II and should use the shared \texttt{\textbackslash{}FanLanCite\{...\}} macro for the primary paper.

\begin{center}\rule{0.5\linewidth}{0.5pt}\end{center}

\subsection{17. Explicit exclusions and chapter boundaries}\label{explicit-exclusions-and-chapter-boundaries}

The chapter will not:

\begin{itemize}
\tightlist
\item
  turn the full BDI literature into a survey beyond the sources selected by Fan \& Lan;
\item
  merge commitment termination policies with reconsideration schedules;
\item
  claim that reflection, logging, or a plan rewrite automatically provides method authority;
\item
  claim that a model proposal is an authorized intention transition;
\item
  claim that CCA solves the full general-purpose B4 lifecycle;
\item
  conflate offline or across-trial method learning with live in-run method authority;
\item
  teach B2 recovery internals, B3 workspace selection, B5 skill verification, or B6 uncertainty/resource governance in full;
\item
  present the proposed Java contracts or architecture as source-prescribed mechanisms;
\item
  invent a universal reconsideration score, fixed trigger threshold, or empirical result that the sources do not report;
\item
  equate hidden-state stability with correctness.
\end{itemize}

\begin{center}\rule{0.5\linewidth}{0.5pt}\end{center}

\subsection{18. Chapter completion and acceptance criteria}\label{chapter-completion-and-acceptance-criteria}

The chapter will be considered ready for author review only when:

\begin{enumerate}
\def\labelenumi{\arabic{enumi}.}
\tightlist
\item
  the historical BDI and MIDCA mechanisms are accurately reconstructed;
\item
  commitment termination and reconsideration scheduling are kept as distinct dimensions;
\item
  all B4 modern examples retain their source-specific boundaries and separate E/D codes;
\item
  the distinction among goal authority, method authority, and model proposal authority is explicit;
\item
  the implementation bridge contains typed contracts, lifecycle state machine, authority matrix, persistence/concurrency strategy, side-effect rules, delegation/suspension semantics, anti-oscillation controls, Java-like interfaces, and a vertical-slice plan;
\item
  the Product X example demonstrates a method switch without goal substitution;
\item
  evaluation preserves Fan \& Lan's required trace, primary outcomes, safety/efficiency measures, three ablations, and falsification criteria;
\item
  cross-bundle boundaries to B1, B2, B3, B5 and B6 are clear;
\item
  citations, glossary, index, seven Mermaid sources, labels, and references compile without unresolved entries;
\item
  the final individual PDF has been compiled with the repository's XeLaTeX/Biber pipeline, rendered to images, visually inspected, and text-extracted;
\item
  the source-fidelity and cross-chapter reviews are recorded;
\item
  no placeholder, mock diagram, or temporary-font build is represented as a final release PDF.
\end{enumerate}

\begin{center}\rule{0.5\linewidth}{0.5pt}\end{center}

\subsection{19. Outline decision and approval record}\label{outline-decision-and-approval-record}

\textbf{Current status:} Proposed outline v1.0; awaiting individual author approval.\\
\textbf{Drafting status:} Not authorized yet. \texttt{Tex/part2/chapter15/chapter.tex} remains a placeholder.\\
\textbf{Next action after approval:} Draft Chapter 15 against this outline, implement the seven planned Mermaid sources, add only necessary shared glossary/index entries, update chapter-specific build and QA checks, compile the individual chapter PDF, inspect all diagrams and pages, and record the Batch 6 checkpoint.

\begin{center}\rule{0.5\linewidth}{0.5pt}\end{center}

\subsection{Primary source}\label{primary-source}

Fan, Haodi, and Zucong Lan. ``From Cognitive Architectures to Language Agents: A Mechanism-Level Review of Lineage, Convergence, and Migration Gaps.'' arXiv:2607.23942 (2026). https://arxiv.org/abs/2607.23942.\autocite{fanlan2026mechanisms}
